\selectlanguage{english}
\noindent\textbf{From cofibrant diagrams to a calculus of homotopy Kan extensions}

The second half of Chapter XVII supplies a construction that the preceding discussions have repeatedly needed: a homotopy theory of diagrams in a category of coefficients, compatible with changing the category of indices. Grothendieck begins with the economical axioms of a category with cofibrations and ends with projective model structures on arbitrary small diagram categories. The intervening work explains why the latter are useful. They provide a uniform source of derived adjunctions, while the earlier constructions on ordered sets explain how the resulting cofibrations can be recognized locally. The text therefore moves between two kinds of argument: a formal calculus of lifting and adjunction, and an inductive construction in which the part of a diagram already built controls the next attachment.

This is one concrete realization of the programme described in \citet[\S69, opening discussion]{Grothendieck1983}. A single derived category does not retain enough information about compatible constructions on complexes or spaces. One must also retain the localized categories of diagrams, together with restriction and the operations adjoint to restriction. The letter to Thomason of 2 April 1991 explains the role of the auxiliary model structures in this programme: the weak equivalences specify the homotopy theory; cofibrations and fibrations make its operations constructible \citep[point~1, printed pp.~3--4]{Grothendieck1991Letter}. Chapter XVII shows both the reach of this principle and the difficulty of implementing it with a minimal axiom system.

A convention must remain visible throughout. In \S3 the diagrams under construction are covariant, written $\mathcal M(I^{\circ})=\operatorname{Fun}(I,\mathcal M)$. In \S\S4--6 the notation returns to $\mathcal M(I)=\operatorname{Fun}(I^{\mathrm{op}},\mathcal M)$. A left Kan extension therefore uses $I/i'$ in the first convention and $(i'\backslash I)^{\mathrm{op}}$ in the second. Likewise, well-founded induction concerns an artinian poset in the covariant convention and a noetherian poset in the presheaf convention. The corrections and commentary preserve the printed text where it slips between these formulations, but do not treat a legitimate change of convention as a transcription error.

\medskip\noindent\textbf{The adjunction must survive the construction}

The first obstruction is already present in ordinary localization. An adjunction $L\dashv R$ does not automatically yield an adjunction between a functor obtained by restricting $L$ to a convenient subcategory and the localization of $R$. An equivalence between the localized subcategory and the whole localization is insufficient. Grothendieck's counterexample in Remark~3.6 is decisive: localizing all arrows in a contractible category can identify it with the terminal category, while selecting a terminal object of the original category does not make that object initial. A left adjoint from the terminal category would have to select an initial object. The defect concerns the universal property of the constructed functor, not its existence as a functor.

Theorem~3.7 repairs the argument for the allowed ordered index categories. A functor is factored through its categorical mapping cylinder, or collage. One part is a closed inclusion; its left extension is extension by the initial object. The other part is a retraction with an initial object in every relevant comma category. Each part admits the localized adjunction already justified by the earlier cofibration arguments. Their composite therefore has the desired adjoint property. This is a proof, not an appeal to a general principle that localization preserves adjunctions without hypotheses.

The subsequent attempt to allow arbitrary targets exposes a separate size issue. Finite colimits do not provide the coproduct indexed by every hom-set of an arbitrary target. For example, finite sets admit the finite-branch cofibration axioms with every map a cofibration and only isomorphisms weak equivalences, but restriction from finite $\mathbb Z$-sets to finite sets has no free-action left adjoint. The commentary gives the argument and states precisely which additional colimits repair the construction. Later results for cocomplete cofibration categories require corresponding infinite coproduct and sequential-colimit axioms; they should not be read back into the finite axioms of these pages.

\medskip\noindent\textbf{Small generation and transfer}

Sections~4 and~5 organize a second route. A special cofibration structure has enough colimits, individually accessible objects, a small family generating its cofibrations in the lifting sense, and the inclusion $C^{\square}\subseteq W$. The accessibility condition is used to run the small object argument. It does not by itself assert that the category has a small generating family of presentable objects; identifying it without qualification with local presentability would strengthen the source's hypotheses.

Given $L\dashv R$, the candidate weak equivalences in the target are detected by $R$, while the candidate cofibrations are built from $L(C)$ by pushouts and transfinite compositions. The central calculation is the lifting identity
\[
 (L(C))^{\square}=R^{-1}(C^{\square}).
\]
To make this into a cofibration structure, the source requires $R$ to preserve the relevant colimits and $RL$ to preserve cofibrations. Preservation of trivial cofibrations is a distinct condition. This distinction explains the difficulty of identifying the transferred fibrations: they are initially defined by lifting against the transferred trivial cofibrations, whose description is the very point at issue.

The free-diagram adjunction is particularly favourable. If $I_0$ is the discrete category of objects of $I$, the left adjoint to restriction along $I_0\to I$ sends a family $(A_j)$ to
\[
 i\longmapsto\coprod_{j\in I} I(i,j)\cdot A_j.
\]
For a map of families, every component of its free-diagram image is a coproduct of coefficient maps. Pushout stability and transfinite closure give the needed coproduct stability. The commentary supplies both the cardinal argument for accessibility of an entire diagram and the elementary transfinite construction of these coproducts.

Section~5 then isolates the remaining model-category problem. Put $TC=C\cap W$, $TF=C^{\square}$ and $F=(TC)^{\square}$. The established factorization gives $W=TF\circ TC$; a retract argument gives $TF=F\cap W$. To obtain a Quillen triple one needs the second factorization, into a trivial cofibration followed by a fibration. A small generating family for $TC$ supplies it. With this second factorization, a pushout and retract argument also proves retract stability of $W$, justifying the later comparison with closed model categories. Under the corresponding preservation hypotheses on $RL$, Proposition~5.8 identifies transferred trivial cofibrations as the lifting closure of $L(TC)$ and transferred fibrations as $R^{-1}(F)$.

The distinction between a generating image, its cellular closure, its retract closure, and its double-orthogonal closure matters here. The equality suggested in footnote~23 cannot be true for the raw image of the free-diagram functor; a worked equivariant example shows why pushouts are indispensable. Proposition~5.1 also contains a displayed equality stronger than the argument preceding the additional hypotheses establishes. The audit records the precise inclusions proved and separates this unresolved strengthening from the later, correctly proved transfer theorem. We do not infer falsity merely from an incomplete justification.

\medskip\noindent\textbf{The local structure of a projective cofibration}

Theorem~6.1 constructs, for every small $I$, the model structure on $\mathcal M(I)$ whose weak equivalences and fibrations are objectwise. Thus every left Kan extension is a left Quillen functor and its derived functor is left adjoint to restriction. Its cofibrations are generally more restrictive than objectwise cofibrations. The rest of the chapter explains what their compatibility condition says.

For a map of arrows, the decisive map is the relative latching map
\[
 X_0\amalg_{X_1}Y_1\longrightarrow Y_0.
\]
Along with the cofibration at the source of the arrow, it characterizes the projective lifting condition. Lemma~6.2.5 states this criterion but leaves its proof unwritten. The commentary proves both directions under exactly the stated pushout hypothesis. No terminal object or model-category factorization is needed for that lemma.

For a well-founded diagram, the same pushout occurs every time one extends an already constructed lift to a new object. Grothendieck uses Zorn's lemma to organize the extensions and writes out the final compatibility cube in detail. The mathematical content of that cube is that one lift from a relative latching object enforces both the pointwise lifting equation and naturality with all previous components. Corollary~6.3 correctly places the noetherian hypothesis on $I$ in the presheaf convention. The contrary word in Theorem~6.1(d) is a genuine printed variance slip: the inverse sequence of nonempty tails of $\mathbb N$, with empty inverse limit, gives a counterexample to the literal artinian version. Chapter XIX, Proposition~3.2 and \S3.4, subsequently uses the corrected noetherian formulation \citep{Grothendieck1990XIX}.

The final proposition identifies restrictions that preserve projective cofibrations. When the connected components of each $I/i'$ have terminal objects, right Kan extension is a product of evaluations. Right lifting classes are closed under products, and adjunction gives the preservation assertion. The commentary also develops a converse test with set-valued presheaves on a poset, making precise how much the terminal-component condition can detect.

\medskip\noindent\textbf{What is settled, and what the commentary adds}

The arbitrary-shape derivator construction is no longer an open problem for model categories with small limits and colimits. Cisinski proves the required homotopy Kan extension theorem without the special-generation restriction \citep[Theorem~6.11 and its dual]{Cisinski2003}; Groth gives a projective--injective treatment in the combinatorial case \citep[Proposition~1.30 and Theorem~1.32]{Groth2013}. For a combinatorial simplicial presentation, rectification identifies the localized diagram categories with $h\operatorname{Fun}(N(I^{\mathrm{op}}),\mathcal C_\infty)$ \citep[Proposition~4.2.4.4]{Lurie2009HTT}. These results concern coherent diagrams, retaining the compatibility which the latching construction makes explicit.

The developments marked \ours\ are proved in the accompanying blocks and make no claim of priority. In particular, a special closed model structure on abelian groups detected by $\operatorname{Hom}(\mathbb Z[1/2],-)$ shows that weak equivalences need not be closed under all filtered colimits. The counterexample addresses the expectation expressed on manuscript p.~97 while respecting the chapter's generation and accessibility hypotheses. Together with the poset restriction criterion and the complete missing arrow-category proof, it turns the commentary into a mathematical continuation of the arguments at the points where a reader needs one.

\medskip\noindent\emph{Commentary and audit AI-produced by GPT-6 Astra (OpenAI, ChatGPT Pro mode), checked and integrated by Claude (Anthropic).}
\nocite{*}
